\chapter[Algebra of the UOD]{Algebra of the Universal Optimality Defect}\label{ch:sec:algebra-of-the-universal-optimality-defect}
\section{Introduction}\label{ch:sec:introduction-analytic_uod-2}
Part~\ref{ch:part:posterior} constructed the posterior manifold together with the
canonical logarithmic coordinate map
$F = Q\circ L, \qquad L(P)=\log P,$
where (Q) denotes the quotient projection removing the additive
logarithmic gauge symmetry. The Universal Optimality Defect (UOD) is the
intrinsic squared distance induced by this geometry:
$\mathcal{D}(P,Q)=\|F(P)-F(Q)\|^{2}.$
The purpose of this chapter is to derive an exact closed formula for
\(\mathcal{D}\) and establish its fundamental algebraic properties.
All later comparison theorems will be derived from this identity.
\section{Canonical Hyperplane Representation}\label{ch:sec:canonical-hyperplane-representation}
Part~\ref{ch:part:posterior} defines (Q) abstractly as the quotient map
\[ Q:\mathbb{R}^{n}\rightarrow
\mathbb{R}^{n}/\operatorname{span}{\mathbf1}. \]
For explicit computations we identify this quotient with the Euclidean
hyperplane
\[ H=\Bigl\{x\in\mathbb{R}^{n}:\sum_{i}
x_i=0\Bigr\}. \]
\begin{lemma}
\label{ch:lem:21-1}
The quotient space
\(\mathbb{R}^{n}/\operatorname{span}{\mathbf1}\) is
canonically isometric to (H). Under this identification,
\[ Q(x)=x-\bar x\,\mathbf1, \qquad
\bar x=\frac1n\sum_{i=1}^{n} x_i. \]
\begin{proof}
Every equivalence class
\(x+\operatorname{span}{\mathbf1}\) contains a unique
vector with zero arithmetic mean, namely
\(x-\bar x\,\mathbf1\). Uniqueness follows because adding a
constant changes the mean by that constant. Hence every quotient class
has a unique representative in (H), defining a linear isomorphism. The
Euclidean norm on (H) therefore induces the canonical quotient norm used
throughout Part~\ref{ch:part:posterior}.
\end{proof}
\end{lemma}
\section{Closed Formula}\label{ch:sec:closed-formula}
\begin{theorem}
\label{ch:thm:21-2}
Let
\[ a_i=\log\frac{P_i}{Q_i}, \qquad
\bar a=\frac1n\sum_{i} a_i. \]
Then
\[ \boxed{
\mathcal{D}(P,Q)
=
\sum_{i=1}^{n}(a_i-\bar a)^2.
} \]
Equivalently,
\[ \boxed{
\mathcal{D}(P,Q)
=
\sum_i a_i^2
-
\frac1n\left(\sum_i a_i\right)^2.
} \]
\begin{proof}
By definition,
$\mathcal{D}(P,Q) = \|Q(\log P-\log Q)\|^{2}.$
Set \(a=\log P-\log Q\). By Lemma~\ref{ch:lem:21-1},
$Q(a)=a-\bar a,\mathbf1.$
Taking the Euclidean norm gives
\[ \mathcal{D}(P,Q) = \|a-\bar a,\mathbf1\|^{2},
\]
whose expansion is
\[ \sum_{i}(a_i-\bar a)^{2} = \sum_{i}
a_i^{2}-\frac1n\left(\sum_{i}
a_i\right)^{2}. \]
\end{proof}
\end{theorem}
\section{Equivalent Forms}\label{ch:sec:equivalent-forms}
\begin{corollary}
\label{ch:cor:21-3}
The Universal Optimality Defect admits the variance representation
\[ \boxed{
\mathcal{D}(P,Q)
=
n\,
\operatorname{Var}_{U_n}
\!\left(
\log\frac{P_i}{Q_i}
\right),
} \]
where \(\operatorname{Var}_{U_n}\) denotes variance with respect
to the uniform distribution on the coordinate index set.
\end{corollary}
\section{Fundamental Properties}\label{ch:sec:fundamental-properties}
\begin{proposition}
\label{ch:prop:21-4}
For all posterior distributions (P,Q):
\begin{itemize}
  \item \(\mathcal{D}(P,Q)\ge0\);
  \item \(\mathcal{D}(P,Q)=0\iff P=Q\);
  \item \(\mathcal{D}(P,Q)=\mathcal{D}(Q,P)\);
  \item \(\mathcal{D}\) is invariant under addition of a common constant in logarithmic coordinates.
\end{itemize}
\begin{proof}
Items (1) and (3) follow from the squared Euclidean norm.
If \(\mathcal{D}(P,Q)=0\), then all centered logarithmic ratios
vanish, so
$\log\frac{P_i}{Q_i}=c$
for every (i). Hence \(P_i=e^{c}Q_i\). Since both distributions sum to one,
$1=\sum_{i}P_i=e^{c}\sum_{i}Q_i=e^{c},$
thus (c=0), proving (P=Q).
Item (4) follows because (Q) annihilates constant vectors.
\end{proof}
\end{proposition}
\section{Interpretation}\label{ch:sec:interpretation}
The UOD is the intrinsic quadratic energy of the logarithmic
probability-ratio vector after quotienting the additive gauge symmetry.
Unlike Kullback--Leibler divergence, which averages logarithmic ratios
with respect to a probability measure, the UOD measures their Euclidean
dispersion in the canonical quotient geometry determined by Part~\ref{ch:part:posterior}.
\section{Outlook}\label{ch:sec:outlook}
Part~\ref{ch:part:variational} develops the mathematical theory of the UOD:
\begin{itemize}
  \item differential structure;\label{ch:sec:differential-uod}
  \item local Riemannian approximation;
  \item functional-analytic properties;\label{ch:sec:functional-uod}
  \item global comparison theorems with Total Variation, Hellinger, $\chi^2$ and
\end{itemize}
    Kullback--Leibler divergence;
\begin{itemize}
  \item applications to entropy and security functionals.
\end{itemize}


\paragraph{Running example: geo-indistinguishability.}
The variance representation $\mathcal D(P,Q)=n\operatorname{Var}_{U_n}(a)$ (Corollary~\ref{ch:cor:21-3}) is particularly illuminating for geo-indistinguishability (Section~\ref{ch:sec:geo-indistinguishability-and-location-privacy}): the log-likelihood ratio $a(x,y)=\log(P(x)/Q(y))$ becomes the Euclidean distance between two points on the plane, and the UOD is proportional to the squared spatial distance---establishing a direct link between geometric privacy and Euclidean geometry.

